\originalchapter{手性场、中心荷与玻色化}
\label{ch:chiral}
\source{8.513相互作用玻色子组后半、作业6--8。}
\originalsection{低能场与热态计数}
一维线性色散把右行与左行激发分开。单支右行玻色模$k_n=2\pi n/L$，$n>0$，激发能$\sum_n(2\pi\hbar v/L)nN_n$。给定总整数能量$n$的简并数是整数分拆数$p(n)$，生成函数$\prod_{m\ge1}(1-q^m)^{-1}$。高温于有限尺寸能级间隔而低温于微观带宽时，
\begin{equation}
\log Z=-\sum_{m>0}\log(1-\e^{-\beta 2\pi\hbar vm/L})\simeq\frac L{2\pi\beta\hbar v}\int_0^\infty[-\log(1-\e^{-x})]\dd x=\frac{\pi L}{12\beta\hbar v}.
\end{equation}
展开$-\log(1-\e^{-x})=\sum_{r\ge1}\e^{-rx}/r$并逐项积分给$\sum r^{-2}=\pi^2/6$。于是$E=\pi L\kb^2T^2/(12\hbar v)$、$C=\pi L\kb^2T/(6\hbar v)$。一支Majorana费米子只有半个复费米自由度，等价中心荷$1/2$，热能系数为以上一半。

采用$ c_{\rm tot}=c_R+c_L$时，一般同速CFT有$E/L=\pi c_{\rm tot}\kb^2T^2/(12\hbar v)$。非手性教科书若用$c=c_R=c_L$，同一式写$\pi c\kb^2T^2/(6\hbar v)$。先说明中心荷计数，才能避免两倍差。手性净值$c_-=c_R-c_L$控制净能量流与手性异常，而热态总密度取决于$c_{\rm tot}$，两者不同。
\originalsection{态密度与虚时间相关}
由正则热能关系求微正则熵。令$E=aT^2$，$S=\int C\dd T/T=2aT=2\sqrt{aE}$，$a=\pi c_{\rm tot}L\kb^2/(12\hbar v)$。大能量密度但仍在低能场论窗口时，
\begin{equation}
D_{\rm many}(E)\sim\exp\left[\sqrt{\frac{\pi c_{\rm tot}LE}{3\hbar v}}\right],
\end{equation}
指数前的幂因子依具体边界扇区与有限尺寸。分拆数渐近$p(n)\sim\e^{\pi\sqrt{2n/3}}/(4\sqrt3n)$代入$E=2\pi\hbar vn/L$也得到$c_{\rm tot}=1$的结果。

反周期Majorana场写$\lambda(x)=L^{-1/2}\sum_k\lambda_k\e^{\ii kx}$，$k=2\pi(n+1/2)/L$，$\lambda_{-k}=\lambda_k^\dagger$。零温$k>0$为湮灭模，所以$\tau>0$的相关函数是$L^{-1}\sum_{k>0}\e^{\ii kx-vk\tau}$。几何级数给
\begin{equation}
G(x,\tau)=\frac{\ii}{2L\sin[\pi(x+\ii v\tau)/L]}\ \longrightarrow\ \frac{\ii}{2\pi(x+\ii v\tau)}.
\end{equation}
$\tau<0$需按费米时间排序换号，并选择相应解析支。局域反对易关系由完整正负动量和得到$\{\lambda(x),\lambda(y)\}=\delta(x-y)$，不能仅以正动量构造实场。
\originalsection{手性玻色场与顶角算符}
紧致场$\phi\sim\phi+2\pi$的多分量作用量为
\begin{equation}
L=\frac\hbar{4\pi}\int\dd x[K_{IJ}\partial_x\phi_I\partial_t\phi_J-V_{IJ}\partial_x\phi_I\partial_x\phi_J].
\end{equation}
$K$为非退化对称整数矩阵，$V$正定确保能量稳定。由于一阶时间导数，共轭动量受约束，正确的Dirac括号给$[\phi_I(x),\phi_J(y)]=\ii\pi(K^{-1})_{IJ}\operatorname{sgn}(x-y)$，整体方向随手性约定。允许单值算符$V_l=\e^{\ii l_I\phi_I}$，$l\in\mathbb Z^N$。BCH公式给远距交换相位$\exp[-\ii\pi l^TK^{-1}l'\operatorname{sgn}(x-y)]$。

$K=(1)$时$\e^{\ii\phi}$具有费米交换符号，相关按$1/(x+\ii v\tau)$，所以可将其识别为手性费米场。左右两支$K=\operatorname{diag}(1,-1)$及密度--密度相互作用可映射一维相互作用费米子。它们的低能关联通常有相互作用依赖的幂律，单粒子准粒子极点消失；这一维Luttinger液体与三维Landau Fermi液体的用语需区分。
\originalsection{分数化与边缘异常}
顶角$l$的拓扑自旋为$\theta_l=\e^{\ii\pi l^TK^{-1}l}$，两种顶角完整绕行相位$M_{ll'}=\e^{2\pi\ii l^TK^{-1}l'}$。局域物理粒子为$l=Kn$；在玻色体系$K$对角须为偶数，使这些局域粒子自旋为1。一般整数$l$不必局域，可作为二维体任意子的边缘终点。二维Abel拓扑序按$l\sim l+Kn$分类，因此种类数$|\det K|$。

单支手性场无法作为具有相同局域自由度的独立一维晶格低能谱：一个有界周期一维能带的右行与左行交叉整体相抵，不能仅有净手性模式。二维体的边缘允许净手性，因为外场下体响应提供异常流入。异常说明边缘守恒方程需与体一起解释，不意味着整个物理体系不守恒。$c_-\ne0$的手性能量流也阻止在保持相应结构的条件下把边缘全部开隙。
